\section{Baseline: Standard GAT with Contrastive Loss}
\label{sec_meth_baseline}

The baseline used for all later comparisons is a standard graph attention network
trained on the synthetic dataset using a contrastive loss -
at inference, the output embeddings are L2-normalised and grouped using $k$-means.
This gives a relatively simple starting point
and sets the bar that any change to the embedding network or grouping method must clear.

The baseline also fixes the main training and evaluation procedure used throughout the experiments -
unless stated otherwise,
later methods inherit the same optimiser, training schedule, and clustering protocol.
This keeps the comparisons consistent
and means that improvements can be linked more directly
to the architectural or methodological change being tested,
rather than to changes in the surrounding training setup.

\subsection{Architecture}
\label{sec_meth_baseline_arch}

The baseline network takes the graph $\mathcal{G}$ constructed in Section~\ref{sec_meth_problem_graph}
and produces an embedding $\mathbf{z}_i \in \mathbb{R}^{d}$ for each node, with $d=128$.
The input to the first GAT layer is the representation $\mathbf{h}_i^{(0)}$
defined in Equation~\eqref{eq:layer0_input} -
in the default no-depth configuration,
this combines the three preprocessed features $\mathbf{x}_i \in \mathbb{R}^{3}$
from Equation~\eqref{eq:node_feature_no_depth}
with a learned 32-dimensional joint-type embedding.
The resulting input dimension is therefore $d_0=35$ -
for the earlier depth-enabled diagnostic,
the four features of Equation~\eqref{eq:node_feature_with_depth} give $d_0=36$.

Message passing is performed by two GATv2 layers~\cite{brody2022gatv2},
each using four attention heads and a hidden dimension of 64.
GATv2 is used instead of the original GAT~\cite{velickovic2018gat}
because its attention mechanism allows the ranking of neighbouring nodes to depend on the query node -
as discussed in Section~\ref{sec_meth_gat},
the original GAT is limited by its static attention formulation,
while GATv2 changes the order of the linear projection and nonlinearity
to produce dynamic attention~\cite{brody2022gatv2}.

This distinction is useful for pose grouping
because not every neighbour should contribute in the same way -
a keypoint may need to treat a nearby joint of the same type
differently from an anatomically related joint of another type.
The importance of a neighbour therefore depends on both nodes involved in the interaction -
for this reason,
GATv2 is kept as the attention primitive throughout the methods developed in this chapter.

After the two message-passing layers,
a linear projection maps each node representation to the 128-dimensional embedding space -
the output is then L2-normalised,
placing each embedding on the unit hypersphere $\mathbb{S}^{d-1}$.
This matches the cosine-based contrastive objective used during training
and also gives the downstream $k$-means stage a bounded representation space.
The complete baseline architecture is summarised in Table~\ref{tab:baseline_arch}.

\begin{table}[H]
\centering
\footnotesize
\begin{tabular}{ll}
\toprule
Component & Setting \\
\midrule
Attention primitive            & GATv2~\cite{brody2022gatv2} \\
Number of GAT layers           & 2 \\
Number of attention heads      & 4 \\
Hidden dimension               & 64 \\
Output (embedding) dimension   & 128 \\
Joint-type embedding dimension & 32 \\
Layer normalisation            & yes \\
Output normalisation           & L2 (unit hypersphere) \\
Dropout                        & 0.1 \\
\bottomrule
\end{tabular}
\caption{Baseline GAT architecture used as the comparison point for the
methods evaluated in this chapter - methods inherit these settings
unless a parameter is explicitly changed or included in a later sweep.
The architecture sweep used to select the final PG-GAT configuration
is described in Section~\ref{sec_meth_pggat_training}.}
\label{tab:baseline_arch}
\end{table}

\subsection{Contrastive loss}
\label{sec_meth_baseline_loss}

The embedding network is trained using a margin-based contrastive loss~\cite{chopra2005contrastive}
defined on cosine similarity.
For a graph containing $N$ keypoints
with ground-truth person labels $\mathbf{y}^\star \in \{1,\dots,K\}^{N}$,
the cosine similarity between two L2-normalised embeddings is
\begin{equation}
    \cos(\mathbf{z}_i,\mathbf{z}_j)
    = \mathbf{z}_i^\top \mathbf{z}_j.
\end{equation}
Training pairs are divided into two sets -
positive pairs contain keypoints belonging to the same person,
\begin{equation}
    \mathcal{P}^{+}
    = \{(i,j) :
    \mathbf{y}^\star_i = \mathbf{y}^\star_j,\; i \neq j\},
\end{equation}
while negative pairs contain keypoints belonging to different people,
\begin{equation}
    \mathcal{P}^{-}
    = \{(i,j) :
    \mathbf{y}^\star_i \neq \mathbf{y}^\star_j\}.
\end{equation}
The contrastive objective is then
\begin{equation}
    \mathcal{L}_{\text{contrastive}}
    =
    \frac{1}{|\mathcal{P}^{+}|}
    \sum_{(i,j) \in \mathcal{P}^{+}}
    \bigl(
        1 - \cos(\mathbf{z}_i,\mathbf{z}_j)
    \bigr)
    \;+\;
    \frac{1}{|\mathcal{P}^{-}|}
    \sum_{(i,j) \in \mathcal{P}^{-}}
    \max\!\bigl(
        0,\;
        \cos(\mathbf{z}_i,\mathbf{z}_j) - \gamma
    \bigr).
    \label{eq:contrastive_loss}
\end{equation}
The first term pulls embeddings belonging to the same person towards a cosine similarity of 1 -
positive pairs keep contributing to the loss however close they already are.
The second term penalises different-person pairs whose similarity exceeds the margin $\gamma$ -
once a negative pair has been pushed below $\gamma$ it no longer contributes,
which keeps the gradient focused on the negative pairs that are still too close together.

The margin is set to $\gamma=0.5$ for the experiments in this work -
this value is later tested as part of the 24-run Bayesian loss sweep
described in Section~\ref{sec_meth_pggat_training},
which confirms it as near-optimal.
The same contrastive formulation is kept across the main embedding experiments,
allowing changes in performance to be attributed more directly to the network architecture
rather than to a different training objective.

The symbol $\gamma$ is used for the contrastive margin throughout this chapter -
this avoids confusion with $m$,
which is already used in Section~\ref{sec_meth_dmon_prelim} to denote the total number of graph edges.

\subsection{Training schedule and grouping read-out}
\label{sec_meth_baseline_training}

The baseline model is trained for 50 epochs
using the 400-scene synthetic training split with a batch size of 4.
Optimisation uses AdamW~\cite{loshchilov2019adamw}
with an initial learning rate of $10^{-3}$ and weight decay of $10^{-4}$ -
a cosine-annealing schedule~\cite{loshchilov2017cosineannealing}
reduces the learning rate to $10^{-5}$ over the course of training.
Gradient norms are clipped at $1.0$ to reduce the effect of occasional large updates,
particularly from batches containing more densely populated scenes.

The model is evaluated every five epochs
on the 100-scene synthetic validation split
using the PGA metric defined in Equation~\eqref{eq:pga} -
rather than using the final training epoch automatically,
the checkpoint with the highest validation PGA is kept for evaluation.

At inference,
the trained network is applied once to each image
to produce the set of node embeddings $\{\mathbf{z}_i\}_{i=1}^{N}$ -
these embeddings are already L2-normalised and are grouped using $k$-means with $k=K$.
The scikit-learn implementation~\cite{scikitlearn} is used
with ten random initialisations and a fixed random seed -
the resulting cluster labels are matched to the ground-truth person labels
using the Hungarian assignment,
after which PGA is calculated using Equation~\eqref{eq:pga}.

Chapter~\ref{chap_4} reports two versions of this baseline -
the first is the depth-enabled configuration
using the four-dimensional node features of Equation~\eqref{eq:node_feature_with_depth}.
This is kept as a diagnostic result because,
as discussed in Section~\ref{sec_meth_data_depth},
exact synthetic depth makes the grouping task unusually easy -
the second is the no-depth configuration using Equation~\eqref{eq:node_feature_no_depth}.
This becomes the main baseline for the experiments that follow.

Unless stated otherwise,
the methods from Section~\ref{sec_meth_learned_heads} onward inherit this no-depth setting
together with the 50-epoch training schedule and optimiser configuration described above -
the architecture settings in Table~\ref{tab:baseline_arch}
are the manually selected baseline values used by the learned-head and SA-DMoN experiments.

PG-GAT later performs a separate Bayesian architecture search,
described in Section~\ref{sec_meth_pggat_training},
which selects a larger three-layer network with a 256-dimensional output -
the two-layer baseline is therefore the bar
that any modification of the embedding network or grouping head must clear at matched compute,
rather than the final architecture used by PG-GAT.